%--- E6: Section 4.2, after the aside "Why G_u and not the full variety" -----
\begin{explained}{the point spectrum and the intrinsic operator $G_u$}
\textbf{Step 1: the decomposition.} By Lemma~\ref{lem:gcd}, $\Ann(u)=G_uJ_u$ with $G_u=\Gcd\Ann(u)$ and $\Var(J_u)$ finite. $G_u$ captures the common polynomial factor of all annihilating equations; $J_u$ captures the residual finite structure. The paper calls $\sigma_p(u)=\Var(J_u)$ the \emph{point spectrum} of $u$: a point $(\tau_0,\xi_0)$ corresponds to the exponential $e^{\tau_0t+\xi_0x}$. ``Spectrum'' is used in an algebraic sense: a finite set of exponent pairs, not a set of Fourier frequencies.

\medskip\textbf{Step 2: why $\Ann(u)=(P)\iff\sigma_p(u)=\emptyset$ and $G_u\sim P$.} Here $G_u\sim P$ means they differ by a nonzero constant.
\begin{itemize}[nosep]
\item Suppose $\sigma_p(u)=\emptyset$, i.e.\ $\Var(J_u)=\emptyset$. Over $\C$ the Nullstellensatz says an ideal without common zeros is the whole ring, so $J_u=\ring$ and $\Ann(u)=G_u\ring=(G_u)$: with empty point spectrum the annihilator is generated by its gcd. If moreover $G_u\sim P$, then $\Ann(u)=(P)$.
\item Conversely if $\Ann(u)=(P)$, the gcd is $P$ up to a constant, the cofactor ideal is $\ring$ and its zero set is empty.
\end{itemize}
Interpretation: the equation is fully identifiable when there is no residual finite point spectrum and the common gcd agrees with the proposed operator.

\medskip\textbf{Step 3: a single exponential.} $u=e^{\tau_0t+\xi_0x}$ has $\partial_tu=\tau_0u$, $\partial_xu=\xi_0u$, so $\Ann(u)=(\tau-\tau_0,\xi-\xi_0)$. The generators have no common nonconstant factor: $G_u=1$, $J_u=(\tau-\tau_0,\xi-\xi_0)$, and
\[
\sigma_p(u)=\{(\tau_0,\xi_0)\}:
\]
one point, exactly the exponential in $u$. Since $G_u=1$ and the spectrum is nonempty, this $u$ does not identify any nonconstant operator as the generator of its full annihilator: $\partial_t-\tau_0$ and $\partial_x-\xi_0$ both annihilate $u$, and neither alone generates $\Ann(u)$.

\medskip\textbf{Step 4: the meaning of $\ring u\cap\EP\cong\ring/J_u$.} $\ring u$ is the span of all derivatives of $u$; $\ring u\cap\EP$ consists of the exponential-polynomials obtainable from $u$ by constant-coefficient operators and linear combinations; the quotient $\ring/J_u$ describes this part algebraically. For the single exponential, $\ring/J_u\cong\C$, matching $\ring u=\spn\{u\}$, one-dimensional and entirely inside $\EP$. In general the finite-dimensional quotient also records \emph{multiplicity} information that the set $\sigma_p(u)$ alone loses.

\medskip\textbf{Step 5: why $G_u$ is the intrinsic operator.} Knowing some $P$ with $P(\partial)u=0$, one wants the operator intrinsic to $u$ rather than $P$ itself, which may contain unnecessary factors. Strategy: (1) factor $P=\prod P_i^{e_i}$; (2) remove factors and test whether the resulting functions are in $\EP$; (3) recover $G_u$ and $J_u$; (4) if $J_u=\ring$, the point spectrum is empty and the annihilator is generated by $G_u$. At most $\sum_ie_i$ membership tests are needed.

\medskip\textbf{Step 6: the repeated-factor example.} The variety $\Var(\Ann u)$ can replace $G_u$ only for reduced $P$, because zero sets do not see multiplicities. The aside above uses $P=\tau^2$ and $u=e^{x^2}$: here $\Ann(u)=(\tau)$, so $G_u=\tau\not\sim P$ while $\Var(\Ann u)=\Var(P)=\{\tau=0\}$. The variety criterion would wrongly say that $u$ identifies $\tau^2$. (An earlier draft used $P=\tau$ for this example, which is reduced and so shows nothing about repeated factors; the revised text uses $\tau^2$.)
\end{explained}
